\begin{table}[!t]
\centering\footnotesize
\setlength{\tabcolsep}{4pt}
\caption{Endpoint Ablations: Variant Minus Full Final Accuracy (Percentage Points)}\label{tab:ablation}
\begin{adjustbox}{max width=\columnwidth}
\begin{tabular}{lrrrr}
\toprule
Variant & Cells & Mean & Min. & Max. \\
\midrule
Fed-MDBSCAN-G, stage 1 only & 12 & $-$0.081 & $-$1.50 & $+$0.84 \\
Geometric-median distance control & 12 & $-$0.081 & $-$1.50 & $+$0.84 \\
Fed-MDBSCAN-G, no activation memory & 12 & 0.000 & 0.00 & 0.00 \\
Fed-MDBSCAN-G, no valve & 12 & $+$0.278 & $-$0.27 & $+$3.03 \\
Trust-region multiplier 1.5 & 12 & $-$0.916 & $-$8.42 & $+$2.90 \\
Trust-region multiplier 2.0 & 12 & $-$0.860 & $-$3.09 & $+$2.95 \\
Trust-region multiplier 3.0 & 12 & $+$0.575 & $-$6.38 & $+$12.82 \\
\bottomrule
\end{tabular}
\end{adjustbox}

\par\vspace{3pt}\parbox{\linewidth}{\footnotesize Matched ablation block. Ranges describe heterogeneity across conditions, not confidence intervals. Removing every layer above the stage-1 acceptance region changes the mean by less than a tenth of a percentage point. Cells are dataset~$\times$~condition means over three seeds; the manuscript makes no equivalence claim from a small mean difference.}
\end{table}
